\documentclass[11pt,a4paper]{article}

\usepackage{amsmath,amssymb,amsthm}
\usepackage{mathtools}
\usepackage{hyperref}
\usepackage{cleveref}
\usepackage[margin=1in]{geometry}
\usepackage{enumitem}
\usepackage{booktabs}

\newtheorem{theorem}{Theorem}[section]
\newtheorem{lemma}[theorem]{Lemma}
\newtheorem{proposition}[theorem]{Proposition}
\newtheorem{corollary}[theorem]{Corollary}
\newtheorem{definition}[theorem]{Definition}
\newtheorem{axiom}{Axiom}
\newtheorem{remark}[theorem]{Remark}
\newtheorem{example}[theorem]{Example}

\newcommand{\Prob}{\mathsf{Prob}}
\newcommand{\ProbProp}{\mathsf{ProbProp}}
\newcommand{\ProbPoset}{\mathsf{ProbPoset}}
\newcommand{\E}{\mathbb{E}}
\newcommand{\pand}{\wedge_p}
\newcommand{\por}{\vee_p}
\newcommand{\pnot}{\neg_p}
\newcommand{\pzero}{\mathbb{0}}
\newcommand{\pone}{\mathbb{1}}
\newcommand{\ple}{\leq_p}
\newcommand{\plt}{<_p}
\newcommand{\padd}{+_p}
\newcommand{\pmul}{\cdot_p}
\newcommand{\psub}{-_p}
\newcommand{\AC}{\mathsf{AC}}
\newcommand{\DC}{\mathsf{DC}}
\newcommand{\CC}{\mathsf{CC}}

\title{Zorn's Lemma via Distributions\\[0.5em]
\large Existence Without Selection}

\author{Probabilistic Foundations Project}

\date{\today}

\begin{document}

\maketitle

\begin{abstract}
Classical Zorn's lemma requires the Axiom of Choice to select a maximal element from chains. We show that probabilistic Zorn's lemma---asserting the existence of \emph{distributions} concentrated on near-maximal elements---avoids this selection problem. The key insight: asserting a measure exists is weaker than asserting a point exists. We currently state probabilistic Zorn as an axiom replacing AC; the conceptual argument for why it should be provable (via finite approximation and Kolmogorov extension) is outlined, with full proof deferred to Paper 5. Classical Zorn is recovered as the $\{0,1\}$ boundary case.
\end{abstract}

\section{Introduction}

\subsection{The Problem with Zorn's Lemma}

Zorn's lemma states: every partially ordered set in which every chain has an upper bound contains a maximal element.

The classical proof requires the Axiom of Choice (AC) or equivalent. Why?

\begin{enumerate}
    \item Consider chains of increasing length
    \item To extend a chain, we must \emph{choose} a strictly greater element
    \item After transfinitely many choices, we reach a maximal element
    \item The selection at each step requires AC
\end{enumerate}

The core issue is \textbf{selection}: we must pick specific elements, not just assert they exist.

\subsection{The Probabilistic Solution}

We replace ``there exists a maximal element'' with ``there exists a distribution concentrated on near-maximal elements.''

Key insight: \textbf{asserting a measure exists is weaker than asserting a point exists}.

\begin{itemize}
    \item A measure can be constructed as a limit of finite approximations
    \item The Kolmogorov extension theorem constructs measures without selection
    \item We never pick a specific maximal element---we construct a probability distribution
\end{itemize}

\subsection{Current Status}

We state probabilistic Zorn as an \textbf{axiom} in our Lean formalization. This axiom is intended to \emph{replace} AC, not supplement it. The conceptual argument for provability is given in \cref{sec:why-no-choice}; a full synthetic proof is the subject of Paper 5.

\section{Background: The Axiom of Choice}

\subsection{Statement}

\begin{axiom}[Axiom of Choice]
For any family of non-empty sets $\{S_i\}_{i \in I}$, there exists a choice function $f : I \to \bigcup_i S_i$ with $f(i) \in S_i$ for all $i$.
\end{axiom}

\subsection{The Hierarchy}

\begin{center}
\begin{tabular}{cl}
\toprule
Axiom & Statement \\
\midrule
$\AC$ & Choice from arbitrary families \\
$\DC$ & Dependent Choice: choice from sequences where each step depends on previous \\
$\CC$ & Countable Choice: choice from countable families \\
\bottomrule
\end{tabular}
\end{center}

Strength: $\AC \Rightarrow \DC \Rightarrow \CC$.

\subsection{Why Zorn Needs AC}

Classical Zorn's proof constructs a transfinite sequence of chain extensions. At limit ordinals, we take unions. At successor ordinals, we must \emph{choose} an extension.

The selection at successor ordinals requires AC (or at least DC for well-orderable posets).

\section{Probabilistic Partial Orders}

\subsection{Definition}

\begin{definition}[Probabilistic Poset]
A \emph{probabilistic partial order} on type $\alpha$ is a structure $P : \ProbPoset(\alpha)$ with:
\[
P.\mathsf{le} : \alpha \to \alpha \to \Prob
\]
where $P.\mathsf{le}(x, y)$ is the probability that $x \leq y$.
\end{definition}

\begin{axiom}[Reflexivity]
$P.\mathsf{le}(x, x) = \pone$.
\end{axiom}

\begin{axiom}[Transitivity]
$P.\mathsf{le}(x, y) \pmul P.\mathsf{le}(y, z) \ple P.\mathsf{le}(x, z)$.
\end{axiom}

\begin{remark}
We do not require antisymmetry at this level; it can be imposed for classical posets.
\end{remark}

\subsection{Strict Ordering}

\begin{definition}[Probabilistic Strict Less-Than]
\[
\plt(x, y) := P.\mathsf{le}(x, y) \pmul (\pone \psub P.\mathsf{le}(y, x))
\]
\end{definition}

This is: ``probability that $x \leq y$ and not $y \leq x$'' = ``probability that $x < y$.''

\begin{theorem}[Irreflexivity]
$\plt(x, x) = \pzero$.
\end{theorem}

\begin{proof}
$\plt(x, x) = P.\mathsf{le}(x, x) \pmul (\pone \psub P.\mathsf{le}(x, x)) = \pone \pmul (\pone \psub \pone) = \pone \pmul \pzero = \pzero$.
\end{proof}

\subsection{Classical Posets}

\begin{definition}[Classical Poset]
A probabilistic poset $P$ is \emph{classical} if $P.\mathsf{le}(x, y) \in \{\pzero, \pone\}$ for all $x, y$.
\end{definition}

This recovers standard partial orders: either $x \leq y$ (probability 1) or not (probability 0).

\section{Near-Maximal Elements}

\subsection{The Key Weakening}

Classically, an element $m$ is \emph{maximal} if: there is no $y$ with $m < y$.

Probabilistically, we weaken this to \emph{$\varepsilon$-near-maximal}: the probability of a strictly greater element is at most $\varepsilon$.

\begin{definition}[Near-Maximal]
An element $x$ is \emph{$\varepsilon$-near-maximal} in $P$ if:
\[
\forall y,\, \plt(x, y) \ple \varepsilon
\]
\end{definition}

\begin{theorem}
Every element is $\pone$-near-maximal (trivially).
\end{theorem}

\begin{proof}
$\plt(x, y) \ple \pone$ for all $x, y$.
\end{proof}

\subsection{Classical Maximality}

\begin{theorem}
For classical $P$: $x$ is $\pzero$-near-maximal iff $x$ is classically maximal.
\end{theorem}

\begin{proof}
$\plt(x, y) \ple \pzero$ forces $\plt(x, y) = \pzero$ (since $\pzero$ is minimal).

For classical $P$: $\plt(x, y) = \pzero$ means $P.\mathsf{le}(x, y) = \pzero$ or $P.\mathsf{le}(y, x) = \pone$.

If $P.\mathsf{le}(x, y) = \pone$, then $P.\mathsf{le}(y, x) = \pone$ (else $\plt(x,y) = \pone \neq \pzero$).

So: $x \leq y$ implies $y \leq x$, i.e., $x$ is maximal.
\end{proof}

\section{Chains and Chain-Completeness}

\subsection{Probabilistic Chains}

\begin{definition}[Probabilistic Chain]
A subset $C \subseteq \alpha$ is a \emph{probabilistic chain} in $P$ if:
\[
\forall x, y \in C,\, P.\mathsf{le}(x, y) \padd P.\mathsf{le}(y, x) = \pone
\]
\end{definition}

This says: for any two elements, the ordering is total (with probability 1, they are comparable).

\subsection{Chain-Completeness}

\begin{definition}[Chain-Complete]
A probabilistic poset $P$ is \emph{chain-complete} if every probabilistic chain has an upper bound:
\[
\forall C,\, (\text{$C$ is a chain}) \Rightarrow \exists\, \mathsf{ub},\, \forall x \in C,\, P.\mathsf{le}(x, \mathsf{ub}) = \pone
\]
\end{definition}

\section{Distributions on Near-Maximals}

\subsection{Probability Distributions}

\begin{definition}[Probability Distribution]
A \emph{probability distribution} on $\alpha$ is a structure $D$ with:
\[
D.\mathsf{weight} : \alpha \to \Prob
\]
satisfying normalization conditions.
\end{definition}

\subsection{Concentrated on Near-Maximals}

\begin{definition}
A distribution $D$ is \emph{concentrated on $\varepsilon$-near-maximals} if:
\[
\forall x,\, D.\mathsf{weight}(x) \neq \pzero \Rightarrow \text{$x$ is $\varepsilon$-near-maximal}
\]
\end{definition}

Every element with positive weight is $\varepsilon$-near-maximal.

\section{Probabilistic Zorn's Lemma}

\subsection{Statement}

\begin{axiom}[Probabilistic Zorn]
Let $P$ be a chain-complete probabilistic poset on a nonempty type $\alpha$. For any $\varepsilon \geq \pzero$, there exists a distribution $D$ on $\alpha$ concentrated on $\varepsilon$-near-maximal elements.
\end{axiom}

\begin{remark}
This is stated as an axiom in our current formalization. It replaces AC.
\end{remark}

\subsection{The Witness Form}

For applications, we often need at least one near-maximal element:

\begin{axiom}[Probabilistic Zorn with Witness]
Under the same conditions, there exists an $\varepsilon$-near-maximal element:
\[
\exists x,\, \forall y,\, \plt(x, y) \ple \varepsilon
\]
\end{axiom}

\subsection{Classical Zorn as Corollary}

\begin{corollary}[Classical Zorn]
For a classical chain-complete poset, there exists a maximal element.
\end{corollary}

\begin{proof}
Apply probabilistic Zorn with $\varepsilon = \pzero$. Get $x$ that is $\pzero$-near-maximal in a classical poset. By \cref{sec:near-maximal-classical}, $x$ is classically maximal.
\end{proof}

\section{Why No Choice is Needed}
\label{sec:why-no-choice}

\subsection{The Conceptual Argument}

The key insight is that we assert \emph{distributions exist}, not \emph{points exist}. Distributions can be constructed without selection:

\begin{enumerate}
    \item \textbf{Finite approximation}: For finite sub-posets, near-maximal elements exist (no choice needed---finite sets have maxima).

    \item \textbf{Projective system}: The distributions on finite sub-posets form a projective system: if $F \subseteq F'$, the distribution on $F'$ projects to the distribution on $F$.

    \item \textbf{Kolmogorov extension}: The projective limit of this system is a distribution on the full poset, concentrated on near-maximals.
\end{enumerate}

\subsection{Where Selection is Avoided}

Classical Zorn requires: ``for each chain, \emph{choose} an element strictly above its supremum.''

Probabilistic Zorn requires: ``the distribution on each finite approximation \emph{exists} (constructively).''

The Kolmogorov extension is constructive given the projective system. No selection is needed.

\subsection{What Remains to Prove}

The full proof requires:
\begin{enumerate}
    \item Formalizing the projective system of finite approximations
    \item Formalizing Kolmogorov extension in our framework
    \item Showing the limit is concentrated on near-maximals
\end{enumerate}

This is deferred to Paper 5, where we develop synthetic probability with enough structure to carry out the construction.

\section{Comparison with Classical Approaches}

\subsection{Solovay's Model}

In ZF + DC + ``all sets of reals are Lebesgue measurable,'' AC fails but many choice consequences hold for measurable sets.

Our approach is orthogonal: we work with probability from the start, not with measurability conditions on classical sets.

\subsection{Reverse Mathematics}

In reverse mathematics, various choice principles are calibrated against subsystems of second-order arithmetic.

Our probabilistic Zorn does not fit neatly into this hierarchy because it asserts measure existence, not point existence.

\subsection{Measurable Selection}

There are classical theorems about measurable selection functions. These require measurability hypotheses and often still use AC.

Our approach avoids selection entirely: we never pick a point, only assert a measure.

\section{Applications}

\subsection{Existence of Bases}

Classically: every vector space has a basis (requires AC).

Probabilistically: every vector space has a distribution concentrated on ``near-bases'' (sets that span to within $\varepsilon$).

\subsection{Existence of Ideals}

Classically: every ring has a maximal ideal (requires AC/Zorn).

Probabilistically: every ring has a distribution concentrated on near-maximal ideals.

\subsection{The Pattern}

Many ``existence of maximal object'' theorems can be reformulated:
\begin{center}
\begin{tabular}{ll}
\toprule
Classical (requires AC) & Probabilistic (no AC) \\
\midrule
$\exists$ maximal $x$ & $\exists$ distribution on $\varepsilon$-near-maximal $x$ \\
$\exists$ basis & $\exists$ distribution on $\varepsilon$-near-bases \\
$\exists$ maximal ideal & $\exists$ distribution on $\varepsilon$-near-maximal ideals \\
\bottomrule
\end{tabular}
\end{center}

The probabilistic versions are weaker but often sufficient for applications (and always sufficient in the limit $\varepsilon \to 0$).

\section{Implementation}

\subsection{Lean 4 Code}

\begin{verbatim}
structure ProbPoset (a : Type) where
  le : a -> a -> Prob

axiom ProbPoset.refl {a : Type} (P : ProbPoset a) (x : a) :
  P.le x x = 1

axiom ProbPoset.trans {a : Type} (P : ProbPoset a) (x y z : a) :
  (P.le x y * P.le y z) <= P.le x z

noncomputable def prob_lt (P : ProbPoset a) (x y : a) : Prob :=
  P.le x y * (1 - P.le y x)

structure NearMaximal (P : ProbPoset a) (x : a) (e : Prob) where
  bound : forall y, prob_lt P x y <= e

structure ChainComplete (P : ProbPoset a) where
  upper_bound : forall C, ProbChain P C ->
                exists ub, forall x, C x -> P.le x ub = 1

axiom prob_zorn {a : Type} [Nonempty a] (P : ProbPoset a) (e : Prob) :
    ChainComplete P -> 0 <= e ->
    exists D : ProbDistribution a, ConcentratedOnNearMaximal P D e

axiom prob_zorn_witness {a : Type} [Nonempty a] (P : ProbPoset a) (e : Prob) :
    ChainComplete P -> 0 <= e ->
    exists x : a, NearMaximal P x e

theorem classical_zorn_from_prob {a : Type} [Nonempty a] (P : ProbPoset a) :
    ClassicalPoset P -> ChainComplete P ->
    exists m, ClassicalMaximal P m := by
  intro hcl hcc
  have h := prob_zorn_witness P 0 hcc (Prob.le_refl 0)
  obtain ⟨x, hnm⟩ := h
  exact ⟨x, fun y hxy => classical_maximal P x hcl hnm y hxy⟩
\end{verbatim}

\subsection{Axiom Count}

Current formalization:
\begin{itemize}
    \item 48 axioms for $\Prob$, conditional expectation, etc.
    \item \texttt{prob\_zorn} and \texttt{prob\_zorn\_witness} as explicit AC replacements
\end{itemize}

Goal (Paper 5): prove \texttt{prob\_zorn} as a theorem.

\section{Discussion}

\subsection{Is This Really ``Without Choice''?}

Strictly speaking, we have \emph{replaced} AC with \texttt{prob\_zorn}. The claim ``no choice needed'' means:

\begin{enumerate}
    \item We do not use classical AC
    \item The probabilistic axiom should be provable from weaker principles
    \item The conceptual argument avoids selection
\end{enumerate}

The full justification requires Paper 5.

\subsection{Why Distributions Matter}

Classical mathematics asks: ``Does $x$ exist?''

Probabilistic mathematics asks: ``What is the distribution over possible $x$?''

The second question is often easier to answer (distributions can be constructed as limits) and often sufficient for applications (we typically use properties that hold with high probability, not properties of specific elements).

\subsection{Relation to Randomized Algorithms}

In computer science, randomized algorithms achieve results ``with high probability'' rather than deterministically. Probabilistic foundations formalize this: existence claims are probabilistic, not absolute.

\section{Conclusion}

Zorn's lemma, classically requiring AC, becomes a theorem about distributions in probabilistic foundations. The key insight: asserting a measure exists is weaker than asserting a point exists, and measures can be constructed without selection.

Currently, probabilistic Zorn is an axiom in our formalization. The conceptual argument for its provability (finite approximation + Kolmogorov extension) is sketched but not yet formalized. Paper 5 will develop the synthetic machinery to complete this proof.

The probabilistic version is the main result. Classical Zorn is a corollary, obtained by restricting to $\{0,1\}$ probabilities and taking $\varepsilon = 0$.

\bibliographystyle{plain}

\end{document}
